\subsection{Mean-field generates monotonically improving policies}

Given that the stochastic perturbations in MC-O-PI are unbiased (\autoref{thm: properties of noise}),
the mean-field update that arises when averaging out these perturbations is
\begin{align}
\label{eq: deterministic mcopi}
    q_{k+1}
    \in
    \big\{ q_k + \lrate_k \update_k ( q_{\pi_k} - q_k )
    :
    \pol_k \in \grpol(q_k),
    \update_k = \update(\start_k, \pol_k, f_k)
    \big\}
\end{align}
with learning rates $(\alpha_k)_{k\ge 0}$, start distributions $(\sigma_k)_{k\ge 0}$ and update functions $(f_k)_{k\ge 0}$.

% If the update distributions $(\update_k)_{k \ge 0}$ are uniform over all actions in each state,
% solutions of \eqref{eq: deterministic mcopi} visit monotonically improving policies and consequently converge to the optimal action values $\qopt$.
% This section first formalises the concept of \textit{policy transition},
% then proves the strict improvement at each transition,
% and finally establishes the convergence of \eqref{eq: deterministic mcopi}.


% deterministic system under uniform updates
% - converges thanks to strong property of "monotone policy improvement"
% - same property as when completely uniform updates, but Tsitsiklis proof breaks
% - can noise cause issues?
% - any counterexample would need noise to make suboptimal point / cycle stable

%%%%%%%%%%%%%%%%%%%%%%%%%%%%%%%%%%%%%%%%%%%%%%%%%%%%%%%%%%%
%%%%%%%%%%%%%%%%%%%%%%%%%%%%%%%%%%%%%%%%%%%%%%%%%%%%%%%%%%%
% \subsubsection{Monotone improvement}


% Using the definitions of greedy policies $\pi \in \grpol(q)$ and $\pi_+ \in \grpol(q_+)$, we can obtain the inequality
% \begin{align}
% % \label{eq: ineq on boundary pol pol}
%     \update(s, \pol) ( q_{\pol}(s, \pol) - q(s, \pol) )
%     \leq \update(s, \pol_+)( q_{\pol}(s, \pol_+) - q(s, \pol) ) .
% \end{align}
% for each state (derivation in appendix).
% Thus, as long as $\update(s, \pol(s)) = \update(s, \pol_+(s))$ for each state, a mean-field update satisfies
% \begin{align}
% % \label{eq: pit in unbiaed lemma for improvement}
%     \qval_\pol(s, \pol) \leq \qval_\pol(s, \pol_+)
% \end{align}
% indicating that $\pol \preceq \pol_+$ by the \hyperref[thm: policy improvement theorem]{Policy Improvement Theorem}.


%%%%%%%%%%%%%%%%%%%%%%%%%%%%%%%%%%%%%%%%%%%%%%%%%%%%%%%%%%%
%%%%%%%%%%%%%%%%%%%%%%%%%%%%%%%%%%%%%%%%%%%%%%%%%%%%%%%%%%%
\subsubsection{Policy improvement at every transition}

Using \autoref{thm: strict PIT},
we show that 
policies generated by the mean-field dynamics in \eqref{eq: deterministic mcopi} improve monotonically
because
\begin{align}
\label{eq: strict improvement theorem}
q_{\pi_{t_n}}(s,\pi_{t_{n+1}}(s)) \ge q_{\pi_{t_n}}(s,\pi_{t_n}(s))
\end{align}
for every state and at every transition.
% at every transition time,

\begin{proposition}
\label{thm: improving policies}
Let \((\pi_k)_{k\ge 0}\) be a sequence of policies generated by the mean-field dynamics in \eqref{eq: deterministic mcopi}
with uniform update distributions over all actions in each state (\autoref{asp: unbiased over policies}).
Let \((t_n)_{n\ge 0}\) be the associated transition times (\autoref{def: transition times}).
Under \autoref{asp: standing},
for every \(n \in \mathbb Z_{\ge 0}\) with \(t_{n+1}<\infty\), policy
\(\pi_{t_{n+1}}\) is better than \(\pi_{t_n}\).
\end{proposition}
% 
\begin{proof}
Let $(q_k)_{k \ge 0}$ be a solution of \eqref{eq: deterministic mcopi} generated with update distributions $(\mu_k)_{k \ge 0}$, and let $(\pi_k)_{k \ge 0}$ be the resulting policy sequence with associated transition times $(t_n)_{n \ge 0}$.


Given policy $\pi_{k+1}$ is greedy with respect to $q_{k+1}$ for every time $k \ge 0$, we have $q_{k+1}(s,\pi_{k+1}(s)) \ge q_{k+1}(s, \pi_k(s))$ for every state $s \in \mathcal S$.
Using \eqref{eq: deterministic mcopi}, this becomes
\begin{multline*}
    q_k(s, \pol_{k+1}(s)) + \lrate_k \update_k(s, \pol_{k+1}(s)) \big( q_{\pi_k}(s, \pol_{k+1}(s)) - q_k(s, \pol_{k+1}(s)) \big)
    \\
    \ge
    q_k(s, \pol_k(s)) + \lrate_k \update_k(s, \pol_k(s)) \big( q_{\pi_k}(s, \pol_k(s)) - q_k(s, \pol_k(s)) \big) .
\end{multline*}
Uniformity over the actions in each state gives
$\mu_k(s,\pi_k(s))=\mu_k(s,\pi_{k+1}(s))=:\mu_k(s)\ge0$.
If $\mu_k(s)=0$, no estimate changes at state $s$. The policy selection rule \eqref{eq:greedy-pol-tie-breaking} therefore retains the previous action, so \eqref{eq: ineq for improvement} holds with equality.
If $\mu_k(s)>0$, rearranging the previous inequality yields
\begin{align*}    
    q_{\pi_k}(s, \pol_{k+1}(s)) - q_{\pi_k}(s, \pol_k(s))
    \ge
    \dfrac{1 - \lrate_k \update_k(s)}{\lrate_k \update_k(s)} \big( q_k(s, \pol_k(s)) - q_k(s, \pol_{k+1}(s)) \big) .
\end{align*}
Since $\pi_k$ is greedy with respect to $q_k$ and, by assumption, $0 < \lrate_k \update_k(s) \le 1$, the right-hand side is nonnegative. As a result, we obtain
\begin{align}
\label{eq: ineq for improvement}
    q_{\pi_k}(s, \pi_{k+1}(s)) \ge q_{\pi_k}(s, \pi_k(s))
\end{align}
for every time $k \ge 0$ and every state $s \in \mathcal S$.

Applying the Policy Improvement Theorem to \eqref{eq: ineq for improvement} gives $v_{\pi_{k+1}}\ge v_{\pi_k}$ at every step, and therefore $q_{\pi_{k+1}}\ge q_{\pi_k}$ as well.
Since the target stays equal to $q_{\pi_{t_n}}$ for $t_n\le k<t_{n+1}$, chaining \eqref{eq: ineq for improvement} at each state also gives \eqref{eq: strict improvement theorem}.
By transitivity, $v_{\pi_{t_{n+1}}}\ge v_{\pi_{t_n}}$ at every finite transition. The two action-value functions differ by definition of a transition time. If their state-value functions were equal, the one-step expectation in the proof of \autoref{thm: strict PIT} would make their action-value functions equal too. Thus the inequality is strict at some state, and $\pi_{t_{n+1}}\succ\pi_{t_n}$.
\end{proof}


\autoref{thm: improving policies} indicates that 
MC-O-PI with uniform updates would visit improving policies
if stochastic perturbations were averaged out, for instance by approximating $q_{\pi_k}$ as the average return from many episodes.

% \begin{proof}
% For every policy $\pol \in \grpol(q)$, $\pol_+ \in \grpol(q_+)$ and state $s \in \sset$ we have
% \begin{align}
%     \label{eq: pol greedy on q}
%     q(s, \pol) &\geq q(s, \pol_+) ,
%     \\
%     \label{eq: pol+ greedy on q+}
%     q_+(s, \pol) &\leq q_+(s, \pol_+) .
% \end{align}
% Assume that $\lrate \update(s,a) < 1$ for every state-action pair (this holds ).
% After multiplying both sides of \eqref{eq: pol greedy on q} with the positive constant $ 1 - \lrate \update(s, \pol_+)$ and subtracting the products from \eqref{eq: pol+ greedy on q+} we obtain
% \begin{multline*}
%     q_+(s, \pol) - ( 1 - \lrate \update(s, \pol_+)) q(s, \pol)
%     \\
%     \leq q_+(s, \pol_+) - ( 1 - \lrate \update(s, \pol_+)) q(s, \pol_+) .
% \end{multline*}
% Further substituting $q_+$ using its definition from \eqref{eq: definition of q+ uniform section} 
% yields
% \begin{align*}
%     \lrate \update(s, \pol)( q_{\pol}(s, \pol) - q(s, \pol) ) + \lrate \update(s, \pol_+) q(s, \pol)
%     \leq \lrate \update(s, \pol_+) q_{\pol}(s, \pol_+) ,
% \end{align*}
% which is equivalent to the inequality
% \begin{align}
% \label{eq: ineq on boundary pol pol}
%     \update(s, \pol) \big( q_{\pol}(s, \pol) - q(s, \pol) \big)
%     \leq \update(s, \pol_+) \big( q_{\pol}(s, \pol_+) - q(s, \pol) \big) .
% \end{align}
% When the update distribution $\update$ is uniform over all actions in each state,
% inequality \eqref{eq: ineq on boundary pol pol} simplifies to
% \begin{align}
% \label{eq: pit in unbiaed lemma for improvement}
%     \qval_\pol(s, \pol) \leq \qval_\pol(s, \pol_+)
% \end{align}
% because $\update(s, \pol) = \update(s, \pol_+)$ for every state.
% As a result, the \hyperref[thm: policy improvement theorem]{Policy Improvement Theorem} indicates that $\pol \preceq \pol_+$.
% \end{proof}

%%%%%%%%%%%%%%%%%%%%%%%%%%%%%%%%%%%%%%%%%%%%%%%%%%%%%%%%%%%
%%%%%%%%%%%%%%%%%%%%%%%%%%%%%%%%%%%%%%%%%%%%%%%%%%%%%%%%%%%
\subsubsection{Global convergence of mean-field solutions}

The monotone policy improvement directly yields global convergence to $\qopt$.

\begin{proposition}
\label{thm: convergence deterministic mcopi}
    Solutions of the mean-field dynamics in \eqref{eq: deterministic mcopi} converge to the optimal action values $\qopt$
    under \autoref{asp: standing},
    if the update distributions are uniform over all actions in each state (\autoref{asp: unbiased over policies}).
\end{proposition}

Technically, this result does not require all conditions in \autoref{asp: standing}. It suffices that $0<\alpha_k\le1$, the policy selection satisfies \eqref{eq:greedy-pol-tie-breaking}, and $\sum_k\alpha_k\mu_k(s)=\infty$ for every state $s$; square summability and \eqref{eq:comparison} are not needed here. We discuss this relaxation in section \ref{sec: relaxing conditions}.


% The mean-field dynamics in Tsitsiklis' setting also exhibit this monotone improvement.
% While their proofs cannot accommodate updates that are only uniform over actions in each state because the crucial commutation step to obtain \eqref{eq: q-Bq decreases} fails,
% the next section confirms that the convergence result nonetheless extends to this broader class of update distributions.


% The following section bridges the gap between the mean-field dynamics and the true stochastic system defined in equation \eqref{eq: mcopi sa}.
% It demonstrates a tight structural coupling between the noisy iterates and a block-wise mean-field solution, suggesting that the stochastic trajectories of MC-O-PI also converge to the optimal action values.

%%%%%%%%%%%%%%%%%%%%%%%%%%%%%%%%%%%%%%%%%%%%%%%%%%%%%%%%%%%%%%%%%%%%%%%
%%%%%%%%%%%%%%%%%%%%%%%%%%%%%%%%%%%%%%%%%%%%%%%%%%%%%%%%%%%%%%%%%%%%%%%
% \subsection{MC-O-PI tracks reinitialised mean-field trajectories}

% Let $(q_k)$ and $(\pi_k)$ be generated by MC-O-PI with learning rates $(\lrate_k)$ and update distributions $(\update_k)$.
% Let $(t_n)$ be the block start times such that $\pi_k = \pi_{t_n}$ for all $k \in [t_n, t_{n+1})$, namely
% \begin{align*}
% t_0 &\coloneqq 0 ,
% \\
% t_{n+1} &\coloneqq \inf\{k>t_n : \pi_k \neq \pi_{t_n}\} .
% \end{align*}
% Define a sequence $(\bar q_k)$ that is reinitialised at the start of each block such that $\bar q_{t_n} = q_{t_n}$, and evolves according to the mean-field update
% \begin{align}
%     \label{eq: deterministic reinitialised mcopi}
%     \bar q_{k+1} &= \bar q_k + \lrate_k \update_k (q_{\pi_{t_n}} - \bar q_k)
% \end{align}
% for all $k \in [t_n, t_{n+1})$.
% Then the stochastic solution $(q_k)$ remains within a shrinking tube around $(\bar q_k)$.

% \begin{lemma}
% \label{thm: shrinking epsilon tube around deterministic solution}
% Under \autoref{asp: standing}, 
% every solution $(q_k)$ of MC-O-PI satisfies almost surely
% \begin{align*}
%     \lim_{n \to \infty} \sup_{t_n \le k < t_{n+1}} \|q_k - \bar q_k \|_\infty = 0 .
% \end{align*}
% \end{lemma}

% \autoref{thm: shrinking epsilon tube around deterministic solution} establishes that MC-O-PI solutions are trapped inside a shrinking tube around the reinitialised mean-field trajectory. 
% Since the mean-field dynamics drive solutions to improving policies,
% we expect that unbiased stochastic perturbations cannot consistently prevent the underlying policy improvement, and hence that MC-O-PI solutions also converge to $\qopt$ when updates are uniform over the actions of each state.

% This monotone improvement in the policy space holds at every step, not only in the limit.
% Thus, such deterministic solutions would never exhibit sliding modes or limit cycles because solutions cannot revisit a policy after leaving its region for a better policy.

% there are several ways to average out noise.
% We use the qualifier \textit{exact} instead of \textit{expected} because the expectation can be taken in several ways.
% Namely, one can either sample an initial state-action pair and filter out the noise by averaging the update over many episodes that start in that state-action pair
% or, additionally, average over all initial state-action pairs, weighted by their respective sampling probability.
% % 
% In the first approach, the start distribution $\sigma_k$ is zero everywhere except for one state-action pair.
% Therefore, the update probabilities in $\update_k$ are conditioned on one particular state-action pair.
% In the second approach, $\sigma_k$ corresponds to the sampling distribution of the initial state-action pair. The update probabilities are then conditioned on that distribution.
% % 
% We will later study the behaviour of exact MC-O-PI for different types of update distributions, regardless of how they are computed.
% We only require that the action values and update distribution are exact at every step, i.e. equal to their expected values.
